\chapter*{Integración de resultados internos LAB-FEN-001 v1}
\markboth{Integración LAB-FEN-001 v1}{Integración LAB-FEN-001 v1}
\addcontentsline{toc}{chapter}{Integración de resultados internos LAB-FEN-001 v1}

\section*{Objeto y estatuto}

\begin{observation}
Esta nota integra en el corpus editorial los resultados ya congelados de LAB-FEN-001-v1. No introduce nuevas métricas, no recalibra umbrales, no modifica trayectorias y no reabre la fase experimental cerrada.
\end{observation}

\begin{observation}
El estatuto de todos los resultados descritos aquí es interno y metodológico. Ninguno constituye por sí mismo validación física externa.
\end{observation}

\section*{Controles y candidata interna}

\begin{itemize}
\item \texttt{LAB-D0-v1} actuó como control negativo de primer orden: la señal histórica conjunta esperada no apareció de forma sistemática.
\item \texttt{LAB-D1-v1} actuó como control positivo rígido de segundo orden y recuperó la señal histórica construida.
\item \texttt{LAB-D2-v1} superó la firma local completa en 8/8 semillas bajo los umbrales congelados.
\end{itemize}

\begin{observation}
La lectura permitida de estos controles es instrumental: el detector separa los casos construidos dentro del contrato declarado. No se infiere de ello una clasificación universal de procesos.
\end{observation}

\section*{Ablaciones}

\begin{observation}
Las ablaciones primarias de \texttt{LAB-D2-v1} fueron diseñadas para separar componentes de la firma observada. La eliminación del pasado inmediato y su sustitución por una señal independiente destruyeron la firma conjunta; la determinización eliminó la aperiodicidad requerida en el control ensayado.
\end{observation}

\begin{observation}
Dentro de LAB-FEN-001-v1, estos resultados son compatibles con una dependencia conjunta de historia causal alineada y elección flexible. Esta lectura está restringida a los mecanismos ensayados y no constituye una necesidad matemática universal ni causalidad física.
\end{observation}

\section*{Atlas paramétrico local}

\begin{observation}
El atlas local predeclarado evaluó un núcleo de 9 puntos y una malla completa de 20 puntos. El núcleo obtuvo 9/9 \texttt{POINT PASS} y la malla observada obtuvo 20/20 \texttt{POINT PASS}.
\end{observation}

\begin{observation}
El resultado permite afirmar que \texttt{LAB-D2-v1} no aparece como punto aislado dentro de las perturbaciones paramétricas v1 declaradas. No permite afirmar que toda extensión paramétrica conserve el régimen ni localiza una frontera exterior.
\end{observation}

\section*{Gate DMO-004}

\begin{observation}
DMO-004 evaluó siete condiciones internas: región paramétrica, acuerdo entre semillas, separación respecto de nulos, criterio explícito de fallo, no-reducibilidad restringida a los límites internos ensayados, regla congelada y trazabilidad completa. El resultado fue 7/7 PASS.
\end{observation}

\begin{observation}
La promoción máxima permitida es \texttt{REGIMEN INTERNO ADMITIDO BAJO LAB-FEN-001-v1 Y DMO-004}. Esta promoción no implica universalidad, irreducibilidad frente a todo modelo posible, interpretación geométrica obligatoria, correspondencia con sistemas naturales ni sustitución de teorías clásicas.
\end{observation}

\begin{observation}
Como referencia externa de terminología, se adopta la noción estándar de cadena de Markov en tiempo discreto \parencite{norris1997markov}. En LAB-FEN-001, los sustitutos Markov-1 se emplean como referencia nula operacional en métricas del protocolo, incluida la evaluación del exceso de dependencia histórica. Esta función debe distinguirse de la definición de \(k^\ast\) en HISTORY-DEPTH-v1: \(k^\ast\) es una clasificación instrumental bajo el protocolo congelado y no se identifica aquí con un orden Markoviano universal.
\end{observation}

\section*{Profundidad histórica operativa}

\begin{observation}
HISTORY-DEPTH-v1 trató la profundidad histórica como variable experimental mediante modelos escalonados. En los cuatro controles construidos se recuperó la escalera operativa: \texttt{LAB-D0} con \(k^\ast=1\), \texttt{LAB-D1} con \(k^\ast=2\), \texttt{LAB-D2} con \(k^\ast=2\) y \texttt{DELAYED} con \(k^\ast=3\).
\end{observation}

\begin{observation}
La lectura permitida para \texttt{LAB-D2-v1} es que la dependencia histórica detectada es compatible con profundidad operativa \(k^\ast=2\) bajo HISTORY-DEPTH-v1. Para el control retardado, la dependencia generativa construida con \(x[t-2]\) es recuperada por el salto correspondiente a profundidad 3.
\end{observation}

\begin{observation}
La ausencia de un hit adicional no demuestra ausencia absoluta de dependencias más profundas. \(k^\ast\) es una profundidad operativa bajo el instrumento congelado, no una constante física ni una geometría.
\end{observation}

\section*{Estabilidad de la profundidad en el atlas}

\begin{observation}
HISTORY-DEPTH-ATLAS-v1 aplicó el mismo instrumento congelado al núcleo paramétrico predeclarado de \texttt{LAB-D2}. Los 9/9 puntos del núcleo conservaron la clasificación operativa \(k^\ast=2\).
\end{observation}

\begin{observation}
Por tanto, puede afirmarse que \(k^\ast=2\) es estable en el núcleo paramétrico predeclarado de \texttt{LAB-D2} bajo HISTORY-DEPTH-v1.
\end{observation}

\begin{observation}
Como observación descriptiva adicional, los 20/20 puntos de la malla v1 y las 160/160 clasificaciones punto-semilla mantuvieron \(k^\ast=2\). Esta cifra no amplía el dominio formal del gate más allá del núcleo predeclarado.
\end{observation}

\section*{Límites explícitos}

\begin{itemize}
\item no se establece que \(k^\ast=2\) sea universal;
\item no se establece que toda familia causal del tipo ensayado tenga \(k^\ast=2\);
\item no se establece estabilidad fuera de la malla estudiada;
\item no se establece una constante física ni memoria física externa;
\item no se establece causalidad física externa;
\item no se impone una geometría a partir de estos resultados;
\item no se presenta MADELEINE como sustitución demostrada de teorías clásicas.
\end{itemize}

\section*{Trazabilidad y cierre de fase}

\begin{observation}
Los resultados aquí resumidos permanecen respaldados por sus especificaciones, ejecutables, resultados estructurados y actas congeladas en la capa de laboratorio. La presente integración editorial no sustituye esos artefactos ni modifica sus valores.
\end{observation}

\begin{observation}
Con HISTORY-DEPTH-ATLAS-v1 se cerró la expansión experimental de LAB-FEN-001 para esta fase de publicación. Las preguntas científicas adicionales quedan como trabajo futuro y requieren protocolos nuevos, separados de la evidencia congelada de esta fase.
\end{observation}
